\section{Vulnerability Type Identification}
\label{sec:problem}
Given a finite set of vulnerability types $T$ and a vulnerable function $f$, the vulnerability type identification (VTI) task associates a subset of types $S \subseteq T$ with the function $f$. \hl{A vulnerable function dataset $D$ consists of $N$ vulnerable functions with their vulnerability types $(f_1, S_1)$, $(f_2, S_2)$,... , $(f_N, S_N)$}. 
%
%
In this work, \textit{we model the problem of vulnerability type identification as multi-label text classification}~\citep{multilabel_cls}. 
%
%
\hl{This is reasonable because each vulnerable function could be considered as a document, and the types of vulnerabilities in the function could be considered as the label/tag set of the document. Note that each function in $D$ is a vulnerable one which could be effectively detected by the existing vulnerability detection techniques}~\cite{ivdetect,linevd,linevul,vuldeeppeaker,sysevr,velvet,poster}.


\subsection{Dataset}
% BigVul: 10,000 vulnerable functions in 13 types.
\textit{Data selecting.} To evaluate VTI approaches, we use BigVul, which is one of the largest vulnerability datasets and provides the vulnerability types of each case. The dataset is collected from 348 real-world C/C++ projects on GitHub, such as Chromium, Linux, Android, PHP, OpenSSL, QEMU, and FFmpeg. The dataset includes about 10,900 vulnerable functions in 13 vulnerability types and 44,603 vulnerable lines of code. 
%
\hl{In this dataset, the types of each vulnerable function are extracted from its corresponding CVE-Details.}

Fig.~\ref{fig:stat_types} and Fig.~\ref{fig:stat_num_labels} show the statistical information of BigVul on vulnerability types and the number of types of vulnerability in each function. Particularly, \textit{Denial of Service (DoS)} and \textit{Overflow} are the two most popular types of vulnerability in the dataset. 
%
This is expected because \textit{Denial of Service (DoS)} and \textit{Overflow} are two of the most frequent vulnerability types~\cite{vul_type_stat}.
%
Meanwhile, \textit{Sql Injection (Sql Inj.)} and \textit{Http Response Splitting (Http R.Spl)} are very rare in BigVul. As seen in Figure~\ref{fig:stat_num_labels}, the vulnerabilities in most of the functions belong to only one type (about 72\%). \hl{Additionally, there are more than 1K vulnerable functions having more than two types.} 
%
%
\begin{figure}
    \centering
    \includegraphics[width=1\columnwidth]{imgs/stat_by_types.pdf}
    \caption{The number of cases by vulnerability types}
    \label{fig:stat_types}
\end{figure}
%
%
\begin{figure}
    \centering
    \includegraphics[width=1\columnwidth]{imgs/stat_by_num_labels.pdf}
    \caption{The number of cases by the number of types}
    \label{fig:stat_num_labels}
\end{figure}


There are some other public vulnerability datasets, but they are not suitable for being used in our experiments. Particularly, SATE IV Juliet~\cite{juliet} is a synthetic vulnerability dataset that is manually created from known vulnerable patterns. However, Chakraborty \etal~\cite{are_we_there} have demonstrated that this dataset contains only simple vulnerabilities which do not reflect real-world vulnerabilities. 
%
%
Meanwhile, Devign~\cite{devign} and Reveal~\cite{are_we_there}, which are constructed from the vulnerabilities in real-world open-source projects, do not provide the vulnerability types for the vulnerable functions. Thus, we do not use these vulnerability datasets in our evaluation experiments.


% Training/testing
\textit{Data splitting.} To evaluate VTI methods, we use a random training/validation/test split ratio of 80:10:10, i.e., the whole BigVul is split into 80\% of training data, 10\% of validating/fine-tuning data, and 10\% of testing data.



\subsection{Evaluation Metrics}
To evaluate VTI approaches, we apply the evaluation metrics which are widely used in multi-label classification studies~\citep{text_multi_label}. For each function $f_i$ in the test set, the types of $f_i$ is represented by a vector $Y_i$ such that for $j \in [1,\|T\|]$, $Y_{ij} = 1$ if $T_j \in T$ is one of the vulnerability types of $f$, otherwise $Y_{ij} = 0$. Let denote vector $Z_i$ the vulnerability type prediction for $f_i$ produced by a VTI technique. For a test set containing $N$ cases (functions), the considering evaluation metrics including \textit{exact match ratio}, \textit{hamming score}, \textit{accuracy}, \textit{macro-average}, \textit{micro-average}, \textit{weighted-average}, \textit{sample-average}, are computed as follows.

\textit{Exact match ratio} indicates the percentage of vulnerable functions (cases) that have \textit{all labels} predicted correctly:
%
$$
\text{\textit{ExactMatchRatio}} = \frac{1}{N}\sum_{i=1}^{N} \mu(Y_i = Z_i)
$$
%
where $\mu$ returns $1$ if $Y_i$ and $Z_i$ are exactly matched, $\forall j \in [1,\|T\|], Y_{ij} = Z_{ij}$, and 0 otherwise. \textit{Exact match ratio} is the most strict one among the metrics. We also use other less strict metrics for the multi-labels classification task. 




\textit{Hamming score} is defined as the proportion of the correctly predicted types to the total number of predicted types and actual types for each case. The overall hamming score is the average across all cases.
%
$$
\text{\textit{HammingScore}} = \frac{1}{N}\sum_{i=1}^{N} \frac{\|Y_i \cap Z_i\|}{\|Y_i \cup Z_i\|}
$$
%
Note that $\|Y_i \cap Z_i\| = \sum_{j=1}^{\|T\|} (Y_{ij} \wedge Z_{ij})$ and $\|Y_i \cup Z_i\| = \sum_{j=1}^{\|T\|} (Y_{ij} \vee Z_{ij})$. 
As seen, \textit{Hamming score} only considers label $1$ in $Y_i$ and $Z_i$. We additionally use \textit{accuracy}, which considers matching for both labels $1$ and $0$.

In this work, \textit{accuracy} is calculated as the following formula:
$$
\text{\textit{Accuracy}} = \frac{1}{N} \sum_{i=1}^{N} \frac{1}{\|T\|} \sum_{j=1}^{\|T\|} \mu(Y_{ij} = Z_{ij})
$$

There are several different methods to measure a multi-label classifier by averaging out the types: \textit{micro-averaging}, \textit{macro-average}, \textit{weighted-average}, and \textit{sample-average}. For \textit{micro-averaging}, all true-positive cases (TPs), true-negative cases (TNs), false-positive cases (FPs), and false-negative cases (FNs) for each type are summed up, and then the average is taken. In the micro-averaging method, we sum up the individual TPs, FPs, and FNs of the system for different sets and then apply them. 

$$
Precision^{micro} = \frac{ \sum_{t \in T} TPs(t) }{ \sum_{t \in T} TPs(t) + FPs(t) }
$$

$$
Recall^{micro} = \frac{ \sum_{t \in T} TPs(t) }{ \sum_{t \in T} TPs(t) + FNs(t) }
$$
And the \textit{micro-average} F1-score will be simply the harmonic mean of the above two equations.
$$
F1^{micro} = \frac{2 \times Precision^{micro} \times Recall^{micro}}{Precision^{micro} + Recall^{micro}}
$$

\textit{Macro-average} is straight forward. We just take the average of the precision and recall of the system on different sets.

$$
Precision^{macro} = \frac{ \sum_{t \in T} Precision(t) }{ \|T\| }
$$

$$
Recall^{macro} = \frac{ \sum_{t \in T} Recall(t) }{ \|T\| }
$$

$$
F1^{macro} = \frac{2 \times Precision^{macro} \times Recall^{macro}}{Precision^{macro} + Recall^{macro}}
$$

\textit{Weighted-average} is simply the average of the precision and recall for individual classes weighted by the support of that class. Meanwhile, to compute \textit{Sample average}, precision, recall, and F1-score are computed for each case and then averaged them.